\documentclass[UTF8,fontset=none,aspectratio=169,10pt]{beamer}
\usepackage{ads-beamer-theme}
\title{高级分布式系统}
\subtitle{第九讲：分布式终止检测}
\author{陈鹏飞\quad 计算机学院}
\institute{中山大学}
\date{2025--2026 学年}

\begin{document}
\begin{frame}[plain]
\begin{tikzpicture}[remember picture,overlay]
  \fill[ADSBlue] (current page.north west) rectangle ([xshift=2.4cm]current page.south west);
  \fill[ADSOrange] ([xshift=2.4cm]current page.north west) rectangle ([xshift=2.55cm]current page.south west);
  \node[anchor=west,align=left] at ([xshift=3.3cm,yshift=1.0cm]current page.west) {\fontsize{32}{38}\selectfont\color{ADSBlue}\bfseries 高级分布式系统};
  \node[anchor=west,align=left] at ([xshift=3.3cm,yshift=-0.4cm]current page.west) {\fontsize{21}{27}\selectfont\color{ADSDarkBlue}第九讲\quad 分布式终止检测};
  \node[anchor=west,align=left] at ([xshift=3.3cm,yshift=-2.2cm]current page.west) {\large\color{ADSGray}所有进程都空闲，真的就意味着整个计算结束了吗？};
  \node[anchor=east] at ([xshift=-0.8cm,yshift=-3.2cm]current page.east) {\small\color{ADSGray}陈鹏飞\quad 中山大学\quad 2025--2026};
  \draw[ADSBlue,line width=0.8pt] ([xshift=3.3cm,yshift=-1.35cm]current page.west) -- ([xshift=12.6cm,yshift=-1.35cm]current page.west);
\end{tikzpicture}
\end{frame}

\sectionframe{一、为什么“大家都闲着”仍然不够？}{分布式计算的终止条件同时涉及进程状态和所有通信信道中的在途消息。}

\begin{frame}{本讲学习目标}
\begin{enumerate}
  \item 精确定义全局终止：局部空闲 + 没有在途消息；
  \item 理解为什么本地观察不能直接推出全局事实；
  \item 掌握快照、权重投掷和生成树波算法的共同思想；
  \item 能说出算法的安全性、活性和消息复杂度；
  \item 把终止检测映射到工作流、训练作业和流处理系统。
\end{enumerate}
\vfill
\callout{主线}{先问“全局谓词是什么”，再问“谁收集足够证据”，最后问“证据在传输途中会不会改变”。}
\end{frame}

\begin{frame}{局部状态与全局状态}
\begin{columns}[T]
\column{.47\linewidth}
\begin{block}{局部可见}
一个进程知道自己的计算是否完成、是否有本地待处理事件。
\end{block}
\begin{block}{全局事实}
所有进程都空闲，并且任意信道都没有尚未到达的基本消息。
\end{block}
\column{.48\linewidth}
\resizebox{.94\linewidth}{!}{%
\begin{tikzpicture}[node distance=7mm and 10mm,>=Latex]
  \nodebox[minimum width=2.1cm]{ADSBlue}{p1t}{idle};
  \nodebox[right=of p1t,minimum width=2.1cm]{ADSOrange}{chan}{in-flight};
  \nodebox[right=of chan,minimum width=2.1cm]{ADSGreen}{p2t}{idle};
  \draw[\arrow,ADSGray] (p1t)--(chan);
  \draw[\arrow,ADSGray] (chan)--(p2t);
\end{tikzpicture}
}
\vfill
\smallnote{两个端点都 idle，但信道中仍有一条消息；此刻不能宣布终止。}
\end{columns}
\end{frame}

\begin{frame}{形式化定义}
设进程 $p_i$ 在时刻 $t$ 的状态为 $p_i(t)$，信道 $c_{ij}$ 中的消息数为 $c_{ij}(t)$。
\[
\text{Terminated}(t_0)\iff
\left(\forall i,\;p_i(t_0)=\text{idle}\right)
\land
\left(\forall i,j,\;c_{ij}(t_0)=0\right).
\]
\vfill
\begin{itemize}
  \item 终止是一个全局谓词，不是某个进程的局部变量；
  \item 进程可以 idle 后再次被消息激活；
  \item 算法必须避免过早宣布，也必须最终宣布。
\end{itemize}
\end{frame}

\begin{frame}{活性与安全性：两个不能混淆的要求}
\begin{tabularx}{\linewidth}{>{\bfseries}p{2.8cm}X}
\toprule
安全性（safety） & 不能产生 false positive：尚有活动进程或在途消息时，不得宣布终止。 \\
活性（liveness） & 如果基础计算确实终止，经过有限时间后必须宣布终止。 \\
非干扰性 & 控制消息不能无限期阻塞基础计算。 \\
信道约束 & 不额外假设新的物理信道，控制信息复用已有通信结构。 \\
\bottomrule
\end{tabularx}
\vfill
\callout{审算法顺序}{先找安全性反例，再证明活性；只证明“最终能停”而不排除过早宣布是不完整的。}
\end{frame}

\begin{frame}{最容易犯的错误}
\begin{itemize}
  \item 错误 1：看到所有进程都 idle 就宣布；遗漏在途消息；
  \item 错误 2：收到了所有“我已经完成”的确认就宣布；确认本身可能早于后续发送；
  \item 错误 3：把一次采样当成全局快照；不同进程的观察时间不一致；
  \item 错误 4：只证明控制消息能到达，不证明活动消息不会绕过检测窗口。
\end{itemize}
\end{frame}

\sectionframe{二、系统模型与推理工具}{算法设计的关键不是“多发几条消息”，而是建立对状态变化和消息在途性的可追踪证据。}

\begin{frame}{终止检测的系统模型}
\begin{itemize}
  \item 每个进程只有两类状态：active 与 idle；
  \item active 进程可以发送基本消息，idle 进程收到消息后被唤醒；
  \item 消息接收可以发生在 active 或 idle 状态；
  \item 发送与接收视为原子动作；
  \item 信道通常可靠、延迟有限，但未必 FIFO。
\end{itemize}
\vfill
\callout{为什么强调模型}{若允许消息丢失、进程崩溃或无限延迟，终止检测要额外结合故障检测和恢复协议。}
\end{frame}

\begin{frame}{因果窗口：一次检测必须覆盖什么？}
\begin{center}
\begin{tikzpicture}[x=1cm,y=1cm,>=Latex]
  \draw[->,ADSGray] (0,0)--(12,0) node[right]{时间};
  \draw[ADSBlue,line width=1.2pt] (1,1)--(4,1) node[right]{P1 active};
  \draw[ADSBlue,line width=1.2pt] (5,1)--(7,1) node[right]{P1 idle};
  \draw[ADSOrange,line width=1.2pt] (2,0.4)--(6,0.4) node[right]{消息在途};
  \draw[ADSGreen,line width=1.2pt] (3,1.8)--(8,1.8) node[right]{P2 active};
  \draw[dashed,ADSGray] (5,0)--(5,2.2);
  \node[font=\scriptsize,anchor=north] at (5,0){检测点};
\end{tikzpicture}
\end{center}
\begin{itemize}
  \item 检测点之前发出的消息可能在检测点之后到达；
  \item 因此必须把“进程状态”和“消息通道状态”一起纳入证据；
  \item 快照算法记录信道状态，权重算法把“未完成性”编码进守恒量。
\end{itemize}
\end{frame}

\begin{frame}{三种算法的共同结构}
\begin{enumerate}
  \item 选择一个可比较的检测轮次、颜色或标签；
  \item 让所有潜在活动源在该轮次中留下证据；
  \item 收集本地状态和在途消息，或守恒的权重总量；
  \item 只有证据闭合，才宣布全局终止。
\end{enumerate}
\vfill
\callout{统一视角}{终止检测本质是“证明没有未来活动来源”，而不是统计当前看起来有多少进程空闲。}
\end{frame}

\sectionframe{三、基于分布式快照的终止检测}{把终止判断转化为一次全局状态记录：所有进程 idle，所有信道为空。}

\begin{frame}{快照算法的直觉}
\begin{itemize}
  \item 当进程从 active 变 idle 时，发起一个带逻辑标签的请求 $R(x,i)$；
  \item 每个进程为自己看到的最大请求拍本地快照；
  \item 请求者收齐所有进程的本地快照后，检查进程与信道状态；
  \item 最大标签对应“最后变 idle 的候选”，能够覆盖更早的请求。
\end{itemize}
\end{frame}

\begin{frame}{逻辑标签 $(x,k)$}
\begin{itemize}
  \item $x$ 表示进程进入 idle 的逻辑轮次；
  \item $k$ 是进程 ID，用于打破相同 $x$ 的并列；
  \item 按字典序比较：$(x,k)>(x',k')$ 当且仅当 $x>x'$，或 $x=x'$ 且 $k>k'$；
  \item 每个进程维护自己见过的最大请求标签；
  \item 最大标签的请求会被所有进程接受，其他请求可能被忽略。
\end{itemize}
\end{frame}

\begin{frame}{快照算法的四条规则}
\scriptsize
\begin{enumerate}
  \item R1：active 进程可以发送基本消息 $B(x)$；
  \item R2：收到 $B(x')$ 时令 $x:=x'+1$；若 idle 则转为 active；
  \item R3：进程进入 idle 时令 $x:=x+1$，设置 $k:=i$，广播 $R(x,k)$ 并为该请求拍快照；
  \item R4：收到 $R(x',k')$ 时，若自己 idle 且标签更大，则更新标签并拍快照；若 active，则至少更新逻辑时间。
\end{enumerate}
\vfill
\callout{关键点}{进程在收到请求时如果已经 active，不能替请求拍“终止快照”，因为它仍可能发送新的基本消息。}
\end{frame}

\begin{frame}{快照判断的安全性}
\begin{itemize}
  \item 若某个进程在请求快照后重新被消息激活，它的本地记录不是 idle；
  \item 若某条基本消息在请求发出前已经发送、但请求快照时仍在途中，信道记录包含该消息；
  \item 因而收齐的全局快照若显示“所有进程 idle 且所有信道为空”，才能宣布终止；
  \item 任何遗漏都会导致 false positive。
\end{itemize}
\end{frame}

\begin{frame}{快照判断的活性}
\begin{itemize}
  \item 基础计算终止后，不再有进程发送新的基本消息；
  \item 最后进入 idle 的进程拥有最大的检测标签；
  \item 该请求最终送达每个进程，所有本地快照都会完成；
  \item 在消息延迟有限的假设下，请求者最终收齐记录并检测到终止。
\end{itemize}
\vfill
\callout{代价}{需要逻辑标签、请求传播和快照收集；它清晰但控制消息可能较多。}
\end{frame}

\sectionframe{四、权重投掷：用守恒量追踪“未完成性”}{把所有活动进程和在途消息携带的未完成性编码成总和为 1 的权重。}

\begin{frame}{权重投掷的核心想法}
\begin{itemize}
  \item 控制代理一开始持有总权重 $1$；
  \item 活动进程和基本消息各持有一个正权重；
  \item 发送消息时，发送者把自己的权重拆成两部分，一部分随消息转移；
  \item 进程变 idle 时，把自己的权重交还控制代理；
  \item 当代理收回权重总和为 $1$，说明没有活动进程、没有基本消息在途。
\end{itemize}
\end{frame}

\begin{frame}{权重投掷的四条规则}
\scriptsize
\begin{enumerate}
  \item R1：发送者将 $W$ 拆为 $W_1+W_2$，自己保留 $W_1$，消息携带 $W_2$；
  \item R2：接收者把消息权重加到本地权重，若 idle 则被激活；
  \item R3：进程转 idle 时发送控制消息 $C(W)$ 给代理，并令本地权重为 0；
  \item R4：代理收到 $C(W)$ 后累加；若代理权重为 1，则宣布终止。
\end{enumerate}
\vfill
\callout{直觉}{权重不会凭空消失；只要活动或消息仍存在，就至少有一个正权重在系统中。}
\end{frame}

\begin{frame}{权重投掷完整例子：1 如何拆分并归还}
\small
控制代理先把权重 1 交给初始活动进程 $P_1$，因此 $W_c=0$。
\begin{center}
\begin{tabularx}{.96\linewidth}{>{\bfseries}p{1.5cm}Xc}
\toprule
步骤 & 动作 & 代理权重 $W_c$ \\
\midrule
1 & $P_1$ 保留 $1/2$，消息 $m_{12}$ 携带 $1/2$。 & 0 \\
2 & $P_2$ 收到 $m_{12}$；保留 $1/4$，消息 $m_{23}$ 携带 $1/4$。 & 0 \\
3 & $P_1$ 变 idle，归还 $1/2$。 & $1/2$ \\
4 & $P_3$ 收到 $m_{23}$，完成后归还 $1/4$。 & $3/4$ \\
5 & $P_2$ 完成并归还自己保留的 $1/4$。 & 1 \\
\bottomrule
\end{tabularx}
\end{center}
\begin{itemize}
  \item 步骤 3 时虽然 $P_1$ 已 idle，但系统仍未终止：$P_2/P_3$ 和消息还持有 $1/2$；
  \item 只有步骤 5 后，所有活动与在途权重都回到代理，才能安全宣布终止。
\end{itemize}
\end{frame}

\begin{frame}{权重算法的两个不变量}
设 $A$ 是活动进程权重集合，$B$ 是基本消息权重集合，$C$ 是控制消息权重集合，代理权重为 $W_c$。
\[
I_1:\quad W_c+\sum_{W\in A\cup B\cup C}W=1
\]
\[
I_2:\quad \forall W\in A\cup B\cup C,\;W>0
\]
\begin{itemize}
  \item $I_1$ 保证总量守恒；
  \item $I_2$ 保证只要还有活动或在途消息，就不可能所有非代理权重为 0。
\end{itemize}
\end{frame}

\begin{frame}{为什么权重为 1 就是安全的？}
\begin{enumerate}
  \item 代理权重达到 1；
  \item 根据 $I_1$，其他所有权重之和为 0；
  \item 根据 $I_2$，集合 $A\cup B\cup C$ 只能为空；
  \item 因而没有活动进程，也没有基本消息或控制消息携带未归还权重；
  \item 于是基础计算已终止，不会过早宣布。
\end{enumerate}
\end{frame}

\begin{frame}{权重算法的工程注意事项}
\begin{itemize}
  \item 正确性依赖基本消息和控制消息可靠送达、恰好计数（或具备去重/重传协议），且进程不在检测期间崩溃；
  \item 浮点数实现不能直接比较“恰好等于 1”，需要有理数、定点数或误差预算；
  \item 权重拆分策略影响数值范围和控制消息大小；
  \item 控制代理可能成为瓶颈或单点故障；
  \item 可以使用树、分区或多个代理扩展，但必须保持守恒关系。
\end{itemize}
\vfill
\callout{理论到工程}{证明中的实数权重是抽象；系统实现要把数值误差、代理故障和消息重试纳入设计。}
\end{frame}

\sectionframe{五、生成树算法：把全局判断分解成子树证据}{树上的局部汇聚降低了通信范围，但必须处理“报告后又被激活”的竞态。}

\begin{frame}{朴素的树上汇聚}
\begin{itemize}
  \item 固定一棵以 $P_0$ 为根的生成树；
  \item 叶子进程 idle 后向父节点发送 token；
  \item 父节点在自己 idle 且收齐所有孩子 token 后，向上发送 token；
  \item 根节点 idle 且收齐孩子 token 时宣布终止。
\end{itemize}
\begin{center}
\begin{tikzpicture}[>=Latex]
  \nodebox{ADSBlue}{r}{P0};
  \nodebox[below left=7mm and 15mm of r]{ADSGreen}{a}{P1};
  \nodebox[below right=7mm and 15mm of r]{ADSGreen}{d}{P2};
  \nodebox[below left=7mm and 7mm of a]{ADSOrange}{b}{P3};
  \nodebox[below right=7mm and 7mm of a]{ADSOrange}{c}{P4};
  \nodebox[below=7mm of d]{ADSOrange}{e}{P5};
  \draw[ADSGray] (r)--(a) (r)--(d) (a)--(b) (a)--(c) (d)--(e);
\end{tikzpicture}
\end{center}
\end{frame}

\begin{frame}{朴素算法的反例}
\begin{enumerate}
  \item 子树中的进程发送 token，父节点认为整个子树空闲；
  \item 之后该进程收到另一个进程发来的消息，重新 active；
  \item 根节点已经收到了旧 token，却不知道子树又产生了活动；
  \item 如果立即宣布终止，就产生 false positive。
\end{enumerate}
\vfill
\callout{真正的难点}{终止检测不是“记录我何时 idle”，而是“证明从检测轮次开始没有新的活动来源”。}
\end{frame}

\begin{frame}{颜色 token：给活动传播留下痕迹}
\begin{itemize}
  \item 初始进程和 token 都是白色；
  \item 进程只要向其他进程发送基本消息，就变黑；
  \item 黑色进程 idle 后发送黑 token，并恢复为白色；
  \item 父节点若收到黑 token，知道子树中发生过跨边消息，必须开启新一轮；
  \item 根节点只有在自己白、自己 idle 且收齐白 token 时才宣布终止。
\end{itemize}
\end{frame}

\begin{frame}{树算法的内向波与外向波}
\begin{center}
\begin{tikzpicture}[node distance=8mm and 11mm,>=Latex]
  \nodebox[minimum width=2.2cm]{ADSBlue}{leaves}{叶子 idle};
  \nodebox[right=of leaves,minimum width=2.2cm]{ADSOrange}{inward}{token 向根汇聚};
  \nodebox[right=of inward,minimum width=2.2cm]{ADSGreen}{root9}{根检查};
  \nodebox[below=of inward,minimum width=2.2cm]{ADSDarkBlue}{repeat}{Repeat 向叶子传播};
  \draw[\arrow,ADSGray] (leaves)--(inward);
  \draw[\arrow,ADSGray] (inward)--(root9);
  \draw[\arrow,ADSGray] (root9)--(repeat);
  \draw[\arrow,ADSGray] (repeat)--(leaves);
\end{tikzpicture}
\end{center}
\begin{itemize}
  \item 白 token 代表本轮没有发现活动传播；黑 token 触发下一轮；
  \item 这是一种把全局检查分解成多轮局部汇聚的思路。
\end{itemize}
\end{frame}

\begin{frame}{消息复杂度}
\begin{itemize}
  \item 最好情况：基础计算早期完成，检测只需一轮树上汇聚，约 $O(N)$ 控制消息；
  \item 最坏情况：每轮都有新的基础消息触发重启，可能达到 $O(NM)$，其中 $M$ 为基础消息数；
  \item 复杂度说明了“简单正确”不等于“消息最优”；
  \item 大规模系统需要增量、分层或与应用状态合并的检测方式。
\end{itemize}
\end{frame}

\sectionframe{六、消息最优与现代工作流映射}{终止检测是分布式工作流、批处理、流式窗口和 AI 训练收敛判断的共同基础。}

\begin{frame}{消息最优算法的出发点}
\begin{itemize}
  \item 固定生成树和领导者可以复用已有拓扑；
  \item 不应让每一条基础消息都触发一轮全树检测；
  \item 目标是把活动传播、计数和检测证据合并，减少额外控制流量；
  \item 若没有领导者，可以先用生成树或 MST 算法选出协调节点。
\end{itemize}
\end{frame}

\begin{frame}{终止检测与快照、共识的关系}
\begin{tabularx}{\linewidth}{>{\bfseries}p{3cm}XX}
\toprule
问题 & 需要收集的证据 & 典型工具 \\
全局终止 & 所有局部 idle + 信道为空 & 快照、权重、树 token \\
全局谓词 & 某个状态谓词是否可能成立 & 一致切、可达性分析 \\
一致决策 & 所有正确进程选择同一值 & 共识、领导者、日志 \\
\bottomrule
\end{tabularx}
\vfill
\callout{区别}{终止检测通常是一次性的全局事实判断；共识要在故障和分歧中持续形成一致决定。}
\end{frame}

\begin{frame}{AI 工作流中的终止条件}
\begin{itemize}
  \item 一个 DAG 的所有上游任务完成，且没有未消费的中间数据；
  \item 所有训练 worker 已完成当前 epoch，并把 checkpoint 写入持久化存储；
  \item 所有推理副本完成预热，流量切换前没有未处理的旧版本请求；
  \item 流式系统的一个窗口关闭，需要同时确认事件时间水位和在途消息。
\end{itemize}
\vfill
\callout{常见误判}{“Pod 都 Completed”仍不等于工作流结束：对象存储写入、队列 ack 和下游触发可能仍在途。}
\end{frame}

\begin{frame}{课堂推演：MapReduce 何时真正结束？}
\begin{enumerate}
  \item 所有 map worker 显示 idle；
  \item shuffle 文件是否全部写入并可见？
  \item 是否还有 reducer 输入在网络或队列中？
  \item reducer 是否可能被迟到的数据重新激活？
  \item 只有满足“任务状态 + 数据传输 + 提交确认”三者闭合，才能宣布 job 完成。
\end{enumerate}
\end{frame}

\begin{frame}{课堂推演：训练作业的安全完成点}
\begin{itemize}
  \item worker 进程退出前是否已发送最后一个梯度并收到聚合确认？
  \item checkpoint 是否已经从本地 NVMe 刷到持久化存储？
  \item 任务控制器是否已记录版本、指标和输出路径？
  \item 任一环节仍有在途消息时，Job 的 Completed 条件都可能过早。
\end{itemize}
\vfill
\callout{设计题}{请给出一个“安全完成”状态机，并说明每条边由谁观测、谁确认。}
\end{frame}

\begin{frame}{算法选择比较}
\begin{tabularx}{\linewidth}{>{\bfseries}p{2.8cm}XXX}
\toprule
算法 & 主要证据 & 优点 & 代价 \\
\midrule
快照式 & 局部状态 + 信道状态 & 语义直接、易证明 & 请求和收集开销 \\
权重式 & 守恒权重 & 精确追踪活动来源 & 代理、数值和故障问题 \\
树 token & 子树汇聚 + 颜色 & 结构简单、易分层 & 多轮重启可能放大消息 \\
\bottomrule
\end{tabularx}
\end{frame}

\sectionframe{七、终止检测完整推演}{从一个在途消息和一次重试开始，判断“看起来完成”是否真的完成。}

\begin{frame}{三进程工作流时间线}
\begin{center}
\begin{tikzpicture}[>=Latex,x=0.16cm,y=1cm]
\draw[ADSGray,->] (0,0)--(68,0); \draw[ADSGray,->] (0,-1.2)--(68,-1.2); \draw[ADSGray,->] (0,-2.4)--(68,-2.4);
\node[anchor=east,font=\scriptsize] at (-1,0){P1};\node[anchor=east,font=\scriptsize] at (-1,-1.2){P2};\node[anchor=east,font=\scriptsize] at (-1,-2.4){P3};
\draw[ADSBlue,very thick] (18,0) circle (2pt) node[above,font=\scriptsize]{send task};
\draw[ADSOrange,very thick] (30,-1.2) circle (2pt) node[above,font=\scriptsize]{active};
\draw[ADSGreen,very thick] (42,-2.4) circle (2pt) node[below,font=\scriptsize]{idle};
\draw[\arrow,ADSBlue] (18,-.1)--(30,-1.1);
\draw[ADSOrange,dashed,->] (30,-1.3)..controls(34,-1.8)and(38,-2.0)..(42,-2.3) node[midway,above,font=\scriptsize]{任务消息};
\end{tikzpicture}
\end{center}
\begin{itemize}
  \item P1 和 P2 可能都显示 idle，但任务消息仍在信道中；
  \item 只有控制证据覆盖应用状态、信道状态和后续依赖，才能宣布终止。
\end{itemize}
\end{frame}

\begin{frame}{快照式检测的四步}
\begin{enumerate}
  \item 检测者发 marker，记录自己的 idle 状态；
  \item 进程收到 marker 后记录本地状态并继续传播；
  \item marker 之间到达的应用消息被记录为信道状态；
  \item 收齐所有输入信道后，检测者检查全局状态与在途消息。
\end{enumerate}
\vfill
\callout{关键}{“所有进程 idle”只能说明节点状态，不包含信道状态。}
\end{frame}

\begin{frame}{权重式检测的守恒直觉}
\begin{itemize}
  \item 控制者持有总权重 1；
  \item 发送任务时，把一部分正权重随消息转移；
  \item 接收并完成任务后，权重返回控制者；
  \item 当所有进程 idle、所有信道为空且权重回到控制者，才满足终止条件；
  \item 浮点数实现必须使用整数 token、精确分数或误差界。
\end{itemize}
\end{frame}

\begin{frame}{黑/白 token 的颜色直觉}
\begin{itemize}
  \item 进程在检测轮次后重新激活，就把 token/自身标成黑色；
  \item 黑色证据沿树向根传播，说明检测期间仍有活动；
  \item 根只有在 token 为白、所有进程 idle 且信道清空时宣布终止；
  \item 颜色机制把“检测期间发生过活动”编码成有限状态。
\end{itemize}
\end{frame}

\begin{frame}{流处理：watermark 不是终止证明}
\begin{enumerate}
  \item watermark 推进到 100，窗口 [0,100) 看起来可以关闭；
  \item 上游重试消息和 checkpoint barrier 仍可能在途；
  \item 下游必须等待数据、状态、存储和控制面确认；
  \item 只有所有依赖闭合，窗口结果才可对外提交。
\end{enumerate}
\end{frame}

\begin{frame}{AI 工作流的完成契约}
\begin{tabularx}{\linewidth}{>{\bfseries}p{3cm}XX}
\toprule
计算 & worker 完成最后一个 step & 不代表 checkpoint 已持久化 \\
通信 & 梯度/控制消息清空 & 不代表对象存储可读 \\
存储 & manifest 原子提交 & 不代表调度器已看到版本 \\
控制面 & Job generation 更新 & 最终完成事件必须幂等 \\
\bottomrule
\end{tabularx}
\end{frame}

\begin{frame}{课堂练习：找出 false positive}
\begin{enumerate}
  \item 三个 worker 都报告 idle；
  \item 控制器收到 Completed 事件；
  \item 之后一个 delayed task 到达 P2，重新触发 checkpoint；
  \item 讨论：原检测漏掉了哪一类状态？如何用 snapshot、weight 或 token 修复？
\end{enumerate}
\end{frame}

\begin{frame}{本讲总结：如何讲清“全局结束”}
\begin{enumerate}
  \item 先把终止定义成全局谓词，而不是本地感觉；
  \item 明确活动进程和在途消息是两类独立的未完成性；
  \item 快照记录状态，权重守恒活动，树 token 汇聚证据；
  \item 安全性防止过早宣布，活性保证最终宣布；
  \item 在 AI 工作流中，Completed 必须覆盖计算、通信、存储和控制面确认。
\end{enumerate}
\vfill
\callout{下一讲衔接}{如果多个进程都想进入临界区，如何在不依赖全局时钟的情况下决定谁先进入？}
\end{frame}

\begin{frame}{终止检测的时序图阅读法}
\begin{itemize}
  \item 先画每个进程的 local state，再画基本消息箭头；
  \item 在每条检测控制消息上标记轮次或颜色；
  \item 找出“看似最后 idle，但之后仍会被激活”的边界；
  \item 检查算法是否记录了该边界之前已经发送、之后才到达的消息。
\end{itemize}
\vfill
\callout{课堂技巧}{遇到复杂图，不要先看算法名字；先找一条可能绕过检测窗口的基本消息。}
\end{frame}

\begin{frame}{安全性证明模板}
\begin{enumerate}
  \item 假设算法宣布终止；
  \item 反设仍存在 active 进程或在途基本消息；
  \item 找出该活动来源携带的控制证据：快照状态、正权重或黑 token；
  \item 证明该证据不可能同时满足“全部 idle/空信道/代理权重为 1”；
  \item 得到矛盾，说明不能 false positive。
\end{enumerate}
\end{frame}

\begin{frame}{活性证明模板}
\begin{enumerate}
  \item 假设基础计算在某时刻 $t_0$ 终止；
  \item 从 $t_0$ 起不再生成新的基本消息；
  \item 证明剩余控制消息在有限延迟下都会送达；
  \item 证明检测状态会单调推进：最大标签、代理权重或白 token 最终闭合；
  \item 因而检测者在有限时间内宣布终止。
\end{enumerate}
\end{frame}

\begin{frame}{控制消息会不会改变基础计算？}
\begin{itemize}
  \item 控制消息不能消耗应用层资源到无限期拖延基本消息；
  \item 控制消息本身若占用同一队列，应有优先级、限流或独立逻辑通道；
  \item 算法安全性不能依赖“控制消息永远先到”，除非模型明确给出优先级；
  \item 工程实现要区分 detection latency 与 application latency。
\end{itemize}
\vfill
\callout{AI 场景}{训练作业的完成确认不能阻塞梯度和 checkpoint 数据面；控制面与数据面要解耦。}
\end{frame}

\begin{frame}{从集中代理到分层终止检测}
\begin{itemize}
  \item 单个控制代理易理解，但在跨区域或大集群上会成为瓶颈；
  \item 可以按机架、可用区或工作流阶段设置局部代理；
  \item 局部代理先判断子系统终止，再由上层汇聚；
  \item 分层算法的关键是证明子系统边界上没有遗留消息或未完成依赖。
\end{itemize}
\end{frame}

\begin{frame}{流处理中的终止：窗口关闭不是源头结束}
\begin{itemize}
  \item 事件时间窗口可以因 watermark 推进而关闭；
  \item 但上游重试、迟到数据和 checkpoint barrier 仍可能在途；
  \item 正确的完成条件需要同时考虑数据时间、处理状态和持久化确认；
  \item 这与快照算法的“进程状态 + 信道状态”结构完全同构。
\end{itemize}
\end{frame}

\begin{frame}{终止检测与幂等性}
\begin{itemize}
  \item 控制器可能重复发送请求、重试收集或重新执行检查；
  \item 如果“宣布完成”不是幂等操作，重复检测会产生重复提交；
  \item 作业系统通常用 generation、epoch、事务 ID 或去重表保护完成事件；
  \item 因而终止检测与分布式提交、exactly-once 语义紧密相连。
\end{itemize}
\end{frame}

\begin{frame}{复杂度讨论：控制消息的三种成本}
\begin{tabularx}{\linewidth}{>{\bfseries}p{3cm}XX}
\toprule
成本 & 问题 & 典型缓解 \\
消息数量 & 大量控制包消耗网络和 CPU & 分层、批量、合并证据 \\
检测延迟 & 过长会拖慢流水线 & 更快的收集路径、局部完成 \\
状态大小 & 标签、权重、token 元数据膨胀 & 摘要、窗口化、压缩 \\
\bottomrule
\end{tabularx}
\end{frame}

\begin{frame}{综合练习：选择算法}
\begin{itemize}
  \item 小型 DAG、节点数量少、需要解释性：优先快照式；
  \item 单一协调服务、要求精确追踪活动来源：可以使用权重式；
  \item 已有稳定生成树、通信拓扑固定：树 token 更经济；
  \item 跨区域 AI 工作流：分层代理 + checkpoint barrier + 幂等完成事件。
\end{itemize}
\vfill
\callout{答题要求}{说明模型假设、检测证据、最坏复杂度和故障后的恢复动作。}
\end{frame}

\begin{frame}{课后题}
\begin{enumerate}
  \item 设计一个三进程快照式终止检测例子，画出一条在途消息导致的 false positive 风险；
  \item 证明权重投掷算法的 $I_1/I_2$ 不变量在发送、接收和转 idle 三个动作下保持；
  \item 分析树 token 算法中“报告后又被激活”的反例，并用颜色机制修复；
  \item 把 Kubernetes Job 的 Completed 条件扩展为“数据、存储、控制器确认”三项检查。
\end{enumerate}
\end{frame}

\begin{frame}{终止检测的适用边界}
\small
\begin{itemize}
  \item 快照式、权重式和树 token 算法都要区分“当前 idle”和“以后不会再产生新活动”；
  \item 权重算法要求发送、接收和转发的权重守恒，工程实现不能用有舍入误差的浮点数直接比较“是否恰好回到 1”；
  \item 若存在丢失、重复、重启或未知的新任务源，必须增加可靠传输、持久化计数器、epoch 或任务注册协议；
  \item Kubernetes Job 的 Completed、流处理窗口关闭和 AI worker 的进程退出，都是局部信号，不能单独推出全局终止。
\end{itemize}
\vfill
\callout{核心判据}{宣布终止前，必须证明应用状态、在途消息、未来活动来源和外部副作用都被同一个检测轮次覆盖。}
\end{frame}

\end{document}
